\chapter{A finite-observation route to the attained minimum}
\label{ch:og-consistency}

The population minimum is an exact statement about ideal process information.
The finite-resolution certificate makes a different promise: a lower bound
from a fixed finite experiment, with its approximation and sampling errors
paid explicitly. This chapter connects the two for the cosine-circle law.
It constructs a sequence of finite experiments whose covered lower bounds
approach the exact minimum. The sufficient rate is deliberately conservative;
neither optimality nor inexpensive laboratory implementation is asserted.

\section{A fixed calibrated family and a sequence of experiments}

Fix a physical reference time $h_0>0$, a calibrated diffusion coefficient
$\gamma>0$, and a calibrated drift bound $\Omega\geq0$. The target process
has the stationary law
\[
 Y_t=\cos\phi_t,\qquad
 d\phi_t=\omega\,dt+\sqrt{2\gamma}\,dW_t\pmod{2\pi},
 \qquad |\omega|\leq\Omega.
\]
The value of $\omega$ is not supplied to the score. Assume exact noiseless
readings for the first theorem. Independent stationary triples, a calibrated
clock, and ordinary hidden-state reversal remain part of the physical
contract. In particular, the triples are not overlapping intervals silently
treated as independent.

Define constants, with their physical time units retained,
\begin{equation}
 M=\gamma+\Omega,\qquad R=\gamma^2+\Omega^2,\qquad
 C=R(2\gamma+\Omega).
 \label{eq:og-consistency-calibration}
\end{equation}
They satisfy $M_{1,+}=M_{1,-}=M$ and $M_{2,+}=M_{2,-}=R$ in the
finite-resolution contract. Indeed,
$\|L^{\vphantom\dagger}f\|_\infty
 =\|L^\dagger f\|_\infty=\sqrt{\gamma^2+\omega^2}\leq M$ and
$\|L^2f\|_\infty=\|(L^\dagger)^2f\|_\infty
 =\gamma^2+\omega^2\leq R$.
The known global diffusivity bound is $\Gamma_f\leq\gamma$.

For an integer $m\geq1$, choose the experiment before its records:
\begin{equation}
 \begin{aligned}
 n_m=m^{12},\qquad h_m=h_0m^{-3},\qquad
 &\varepsilon_m=\gamma m^{-2},\\
 q(y)=-\gamma y,
 \quad a(y)=\gamma(1-y^2),
 &r_m(y)=\mathbf 1_{\{a(y)\geq\varepsilon_m\}}.
 \end{aligned}
 \label{eq:og-consistency-design}
\end{equation}
This is the existing weighted three-time score with envelope floor
$\varepsilon_m$, $Q=\gamma$, and noiseless range $[-1,1]$.
Every excluded center still consumes a complete record. The envelope is
exact under the stated circle calibration; it is not estimated from the test
increments.

Let $W_{m,i}$ be the corrected score in
Equation~\eqref{eq:og-finite-score}, and let $\mu_m=\mathbb E W_{m,i}$.
The correction uses
\[
 \beta_g=h_m R/2,\qquad \beta_K=h_mMR,
 \qquad 2|q|\beta_g+\beta_K\leq Ch_m.
\]
The known-variance Bernstein option in the finite-resolution implementation
applies directly. No new concentration algorithm is required.

\section{An explicit population and confidence bound}

The ideal truncated local target is
\[
 J_m=\sigma\,\mathbb P\{\Gamma_f\geq\varepsilon_m\},
 \qquad \sigma=\omega^2/\gamma.
\]
Uniform stationary phase gives
\begin{equation}
 \sigma-J_m
 =\sigma\frac2\pi\arcsin(1/m)\leq\frac{\sigma}{m}.
 \label{eq:og-consistency-tail}
\end{equation}
The last inequality uses $\arcsin u\leq\pi u/2$ for $0\leq u\leq1$.
Since $q=g$ and $a=\Gamma_f$, the corrected finite-step numerator lies
between its ideal conditional numerator minus twice its bias envelope and
that ideal numerator. Therefore
\begin{equation}
 J_m-\frac{2Ch_m}{\varepsilon_m}
 \leq\mu_m\leq J_m\leq\sigma,
 \qquad
 0\leq\sigma-\mu_m
 \leq\frac{\sigma+2Ch_0/\gamma}{m}.
 \label{eq:og-consistency-population}
\end{equation}
This deliberately uses the general calibrated bias bound rather than
substituting the unknown $\omega$ into an exact finite-step correction.

For explicit concentration constants set
\begin{align}
 H&=\frac{(2\gamma+\Omega+\sqrt{2\gamma/h_0})^2+Ch_0}{\gamma},\\
 B_0&=\frac{h_0^{-2}+(2/h_0+\gamma)^2+Ch_0}{\gamma}.
 \label{eq:og-consistency-constants}
\end{align}
The finite-resolution second-moment bound and range width satisfy
\begin{equation}
 \mathbb E W_{m,i}^2\leq H^2m^{10},\qquad
 B_{m,+}-B_{m,-}\leq B_0m^8.
 \label{eq:og-consistency-moments}
\end{equation}
For example, the square root of the second-moment envelope is at most
\[
 \frac{(\gamma+M+\sqrt{2\gamma/h_m})^2+Ch_m}{\varepsilon_m}
 \leq Hm^5,
\]
because $m\geq1$. The range bound follows from
Equation~\eqref{eq:og-finite-range} with $R=2$ for the detector width;
this detector-width symbol is separate from the calibrated generator
constant $R$ in~\eqref{eq:og-consistency-calibration}.

Choose a desired two-sided failure probability $\alpha\in(0,1)$ and put
$\ell=\log(2/\alpha)$. Choose the calibrated range-width and
second-moment bounds $b_m=B_0m^8$ and $v_m=H^2m^{10}$, or any valid
smaller bounds, and set, in exact real arithmetic,
\begin{equation}
 \rho_m=\frac{b_m\ell}{3n_m}
 +\sqrt{\frac{2v_m\ell}{n_m}
        +\left(\frac{b_m\ell}{3n_m}\right)^2},
 \qquad
 \widehat L_m=[\overline W_m-\rho_m]_+.
 \label{eq:og-consistency-estimator}
\end{equation}
This uses the same known-variance Bernstein construction as the implementation, with
one-sided allocation $\alpha/2$. Applying the same inequality to both signs
and taking a union bound gives
$|\overline W_m-\mu_m|\leq\rho_m$ with probability at least $1-\alpha$.
The deterministic envelopes imply
\begin{equation}
 \rho_m\leq
 \frac{H\sqrt{2\ell}}{m}
 +\frac{2B_0\ell}{3m^4}.
 \label{eq:og-consistency-radius}
\end{equation}

\begin{theorem}[Constructive consistency at the attained circle minimum]
\label{thm:og-consistency}
For the calibrated experiment above, with probability at least $1-\alpha$,
\begin{equation}
 \boxed{\displaystyle
 0\leq\sigma-\widehat L_m
 \leq
 \frac{\sigma+2Ch_0/\gamma+2H\sqrt{2\log(2/\alpha)}}{m}
 +\frac{4B_0\log(2/\alpha)}{3m^4}.}
 \label{eq:og-consistency-main}
\end{equation}
Replacing $\sigma$ on the right by $\Omega^2/\gamma$ gives a completely
calibrated uniform constant for $|\omega|\leq\Omega$.
For every fixed $\alpha$, $\widehat L_m\to\sigma$ in probability as
$m\to\infty$.
\end{theorem}

\begin{proof}
On the two-sided Bernstein event,
$\overline W_m-\rho_m\leq\mu_m\leq\sigma$ and
$\overline W_m-\rho_m\geq\mu_m-2\rho_m$. Taking the positive part
preserves the upper inequality because $\sigma\geq0$, and can only improve
the lower inequality. Equations~\eqref{eq:og-consistency-population} and
\eqref{eq:og-consistency-radius} establish the displayed bound.

Convergence in probability is not deduced merely by fixing a nonzero
failure probability in that bound. Rather,
\[
 \operatorname{Var}(\overline W_m)
 \leq\frac{H^2m^{10}}{m^{12}}=\frac{H^2}{m^2}\longrightarrow0.
\]
Chebyshev's inequality gives
$\overline W_m-\mu_m\to0$ in probability, while
$\mu_m\to\sigma$ and $\rho_m\to0$ deterministically. Continuity of
the positive-part map proves the stated convergence, including $\omega=0$.
\end{proof}

The same result holds with a conservative rational radius that actually
vanishes in an executable implementation. Choose once, by outward
enclosure, rational constants $\overline H\geq H$,
$\overline B\geq B_0$, $\overline\ell\geq\log(2/\alpha)$ and
$\overline k\geq\sqrt{2\overline\ell}$. Use
\begin{equation}
 \overline\rho_m=
 \frac{\overline H\,\overline k}{m}
       +\frac{2\overline B\,\overline\ell}{3m^4}
 \label{eq:og-consistency-rational-radius}
\end{equation}
and exact rational division by the integer powers of $m$. This radius
dominates the exact Bernstein radius, gives the same confidence argument,
and tends to zero; the rate bound uses the corresponding enclosed constants.
The reference wrapper implements this construction without changing the
existing finite-resolution module. A fixed absolute rounding grid applied
directly to the final square root would leave a nonzero numerical floor,
so the literal fixed-precision final-radius calculation is not used to claim
asymptotic consistency.

If almost-sure convergence is desired for a sequence of experiments, use a
summable sequence of two-sided failure budgets, for example
$\alpha_m=m^{-2}$ for $m\geq2$. The bound then has order
$\sqrt{\log m}/m$, and Borel--Cantelli gives eventual success and
almost-sure convergence. Dependence between different experiments does not
invalidate that summability argument; independence of records within each
experiment remains required.

\section{Controlled score perturbations and resource costs}

The same conclusion survives a calibrated uniform perturbation of each
already defined score. Suppose the stored iid scores satisfy
$|\widetilde W_{m,i}-W_{m,i}|\leq\eta_m$, with the same fixed selection,
weight and drift rule. Expand the score range by $\eta_m$ at both ends,
replace the square root of the second-moment bound by
$\sqrt{v_m}+\eta_m$, and subtract $\eta_m$ from the final lower bound,
as in the existing quantization interface. Applying the two-sided argument
to the stored scores gives
\begin{equation}
 0\leq\sigma-\widehat L_m^{\rm stored}
 \leq\sigma-\mu_m+2\widetilde\rho_m+2\eta_m
 \label{eq:og-consistency-perturbed}
\end{equation}
with probability at least $1-\alpha$. In particular, $\eta_m=O(m^{-1})$
preserves the sufficient rate and consistency. This covers a proved
score-rounding envelope. It does not assume that rounding the central
detector value leaves the selector, local diffusivity or weight unchanged.
A finite-resolution detector must justify those changes, including the
boundary of the retention region, before this result applies.

Alternatively, retain the radius computed for the ideal iid scores and
assume a certified deterministic error
$|\widetilde{\overline W}_m-\overline W_m|\leq\eta_m$ in their stored
mean. Subtracting $\eta_m$ then gives the same performance inequality with
the ideal-score radius and the term $2\eta_m$. This follows directly on
the ideal-score concentration event and does not require a fresh independence
claim about numerical errors. The rational consistency wrapper uses this
version. The error envelope must still cover the actual fixed score rule.

For one experiment at level $m$, the declared costs are
\begin{equation}
 \text{stationary preparations}=m^{12},\qquad
 \text{endpoint readings}=3m^{12},\qquad
 T_m=2n_mh_m=2h_0m^9.
 \label{eq:og-consistency-resources}
\end{equation}
Thus the conservative sufficient accuracy rate is $n_m^{-1/12}$ in the
number of preparations, or $(T_m/(2h_0))^{-1/9}$ in declared evolution
time. These are not minimax exponents. They account for increasingly fine
timing, a shrinking geometry floor and growing sample size using broad
uniform bounds. Preparation, reset, sensor bandwidth, calibration,
data handling, and energy costs are additional; $T_m$ is not total
laboratory duration. If every level is run successively, the cumulative
cost is the sum over those levels, rather than the single-level cost above.

The theorem closes a precise existence question: under the stated
calibration and access assumptions, finite covered experiments can approach
the population minimum. It does not establish a practical advantage at
every resource budget, reconstruct the hidden architecture, or transfer the
same rate to arbitrary nonlinear detectors and diffusions.
